% --- Añadido por 2latex: lo que necesita el contenido convertido desde Obsidian. ---
% Se carga en el preámbulo de la plantilla. Si la plantilla ya carga alguno de estos paquetes,
% LaTeX ignora la segunda carga (sin opciones no hay conflicto).
\usepackage{amsmath,amssymb}
\usepackage{graphicx}
\usepackage{longtable,booktabs,array,calc}
\usepackage{xcolor}
\usepackage[most]{tcolorbox}
\usepackage{hyperref}
\hypersetup{hidelinks}

% Macros que usa la salida de pandoc
\providecommand{\tightlist}{\setlength{\itemsep}{0pt}\setlength{\parskip}{0pt}}
\providecommand{\passthrough}[1]{#1}
\providecommand{\st}[1]{#1}
\providecommand{\ul}[1]{\underline{#1}}
\makeatletter
\providecommand{\pandocbounded}[1]{%
  \begingroup\sbox0{#1}%
  \ifdim\wd0>\linewidth\resizebox{\linewidth}{!}{\usebox0}\else\usebox0\fi
  \endgroup}
\makeatother

% Callouts de Obsidian: \begin{callout}{tipo}{título}
\definecolor{callout-note}{HTML}{448AFF}
\definecolor{callout-abstract}{HTML}{00B0FF}
\definecolor{callout-info}{HTML}{00B8D4}
\definecolor{callout-tip}{HTML}{00BFA5}
\definecolor{callout-success}{HTML}{00C853}
\definecolor{callout-question}{HTML}{64DD17}
\definecolor{callout-warning}{HTML}{FF9100}
\definecolor{callout-danger}{HTML}{FF1744}
\definecolor{callout-example}{HTML}{7C4DFF}
\definecolor{callout-quote}{HTML}{9E9E9E}
\newtcolorbox{callout}[2]{%
  enhanced, breakable,
  colframe=callout-#1!80!black, colback=callout-#1!6, coltitle=white,
  colbacktitle=callout-#1!80!black, fonttitle=\bfseries,
  boxrule=0.6pt, arc=2pt, left=6pt, right=6pt, top=4pt, bottom=4pt,
  before skip=8pt, after skip=8pt,
  title={#2}}
% --- Fin de lo añadido por 2latex ---
